%% kapa: exam=1 round=24 year=2013 period=1 form=A subject=civil_law subject_ko=민법 seq=34 number=34
%% source: Q-Net 감정평가사 기출문제 (cst00309) artlSeq=5073802 file=extracted/1008/5073802/2013년 제24회 감정평가사 1차시험 기출문제/1교시(A형).hwp
%% conversion: kordoc 4.24.51 -> tools/gen1.py (auto); review=proofread-sonnet; flags=none
\begin{question}{34}
전세권에 관한 설명으로 옳지 않은 것은? (다툼이 있으면 판례에 의함)
\begin{choices}
  \item 건물 일부의 전세권자는 나머지 건물 부분에 대한 경매신청권이 없다.
  \item 토지전세권의 경우 전세권의 법정갱신이 인정된다.
  \item 토지전세권의 존속기간이 만료한 경우 전세권의 용익물권적 권능은 전세권설정등기의 말소 없이도 당연히 소멸한다.
  \item 전세권설정자가 제3자에게 전세 목적 부동산을 양도한 경우 그는 원칙적으로 전세금반환의무를 면한다.
  \item 전세권자가 장차 전세권이 소멸하여 전세금 반환채권이 발생하는 것을 조건으로 그 장래의 조건부 채권을 제3자에게 양도한 경우 그 양도는 유효하다.
\end{choices}
\end{question}
\answer{2}
